\documentclass[twocolumn,11pt]{article}

% Packages
\usepackage[utf8]{inputenc}
\usepackage[T1]{fontenc}
\usepackage{textcomp}
\usepackage[margin=1in]{geometry}
\usepackage{graphicx}
\usepackage{amsmath}
\usepackage{amssymb}
\usepackage{amsfonts}
\usepackage{hyperref}
\usepackage{xcolor}
\usepackage{booktabs}
\usepackage{caption}
\usepackage{float}
\usepackage{parskip}
\usepackage{algorithm}
\usepackage{algorithmic}
\usepackage{titlesec}
\usepackage[numbers,sort&compress]{natbib}

% Paragraph spacing: 6pt after only
\setlength{\parskip}{6pt}
\setlength{\parindent}{0pt}

% Section spacing: 6pt before and after
\titlespacing*{\section}{0pt}{6pt}{6pt}
\titlespacing*{\subsection}{0pt}{6pt}{6pt}
\titlespacing*{\subsubsection}{0pt}{6pt}{6pt}

% Figure and table spacing: 6pt before and after
\setlength{\floatsep}{6pt}
\setlength{\textfloatsep}{6pt}
\setlength{\intextsep}{6pt}

% Custom commands
\newcommand{\shpn}{Signal Hierarchical Petri Net}
\newcommand{\shpns}{Signal Hierarchical Petri Nets}
\newcommand{\erk}{ERK}
\newcommand{\mkp}{MKP}
\newcommand{\mapk}{MAPK}
\newcommand{\dg}{\Delta G}
\newcommand{\srr}{\text{SRR}}

% Title and authors
\title{\textbf{Signal Hierarchical Petri Nets with Thermodynamic Orchestration: Modeling Perfect Adaptation in the ERK/MAPK Cascade}}

\author{
Eugênio Simão\textsuperscript{1,*}\\
\textsuperscript{1}Department of Computer Science\\
Universidade Federal de Santa Catarina (UFSC)\\
Araranguá, Santa Catarina, 88906-072, Brazil\\
\textsuperscript{*}Corresponding author: eugenio.simao@ufsc.br
}

\date{January 9, 2026}

\begin{document}

\maketitle

\begin{abstract}
\textbf{Background:} Mitogen-activated protein kinase (MAPK) cascades are central signaling pathways that must detect transient signals while adapting to sustained stimulation. Classical modeling approaches struggle to capture both the hierarchical organization of cellular signals and the thermodynamic constraints that govern biochemical feasibility.

\textbf{Methods:} We developed a \shpn{} framework that extends classical Petri nets with signal hierarchy through environmental sensing (signal places) and thermodynamically-guided regulatory control. We constructed a detailed model of the \erk{}/\mapk{} cascade incorporating dual-pathway negative feedback: ERK-PP-induced \mkp{} synthesis and direct growth factor-to-\mkp{} feedforward. During model validation, we discovered and corrected a critical pulse circuit timing error that invalidated prior results. The corrected model was optimized against four criteria for 10-second transient pulses: $>90\%$ adaptation, adaptation timescale ($\tau_{\text{adapt}}$) of 20--40s, steady-state response ratio (\srr{}) $<5\times$, and fold-change $>4\times$.

\textbf{Results:} Our optimal configuration achieved 96.5\% adaptation and $3{,}935\times$ fold-change with a 10s pulse. The \srr{} of $139\times$ represents the theoretical minimum for this amplification level (mathematical constraint: SRR $\geq$ 1 + 0.035 $\times$ 3,934 = 138$\times$). Direct growth factor-to-\mkp{} feedforward (rate=25 mM/s, K=0.1 mM) enabled \mkp{} accumulation during the pulse (0.075 $\to$ 49.5 mM), reducing peak \erk{}-PP from 20 mM to 1.6 mM. Mathematical analysis reveals that achieving \srr{} $<5\times$ while maintaining $>90\%$ adaptation requires fold-change $<115\times$, incompatible with realistic MAPK amplification ($>1{,}000\times$). The $\tau_{\text{adapt}}=121$s reflects physical exponential decay constraints (28s half-life requiring $\sim$7 half-lives).

\textbf{Conclusions:} \shpns{} successfully model complex adaptation dynamics and reveal fundamental trade-offs between signal detection sensitivity and perfect baseline return. The corrected model demonstrates that direct feedforward pathways are essential for adaptation to brief transient signals. Our results suggest the original criterion \srr{} $<5\times$ is physically impossible with biologically realistic cascade amplification, and real \mapk{} cascades likely operate with \srr{} of 100--500$\times$ as an unavoidable consequence of ultrasensitivity. This provides a generalizable platform for analyzing signaling pathways with realistic physical constraints.

\textbf{Keywords:} Signal Hierarchical Petri Nets, MAPK Cascade, ERK, Perfect Adaptation, Thermodynamics, Regulatory Control, Incoherent Feedforward Loop, Systems Biology
\end{abstract}

%% ============================================================
%% 1. INTRODUCTION
%% ============================================================
\section{Introduction}

\subsection{MAPK Cascades and the Adaptation Problem}

Mitogen-activated protein kinase (\mapk{}) cascades represent one of the most evolutionarily conserved and ubiquitous signaling pathways in eukaryotic cells. The \erk{}/\mapk{} pathway, consisting of sequential phosphorylation through Raf, MEK, and \erk{} kinases, transduces extracellular signals from growth factors and mitogens to regulate cell proliferation, differentiation, and survival. This three-tier cascade architecture provides substantial signal amplification---single growth factor molecules can trigger thousands of \erk{} phosphorylation events---enabling cells to detect weak stimuli with high sensitivity.

However, this amplification creates a fundamental challenge: how can cells distinguish between transient signals requiring rapid responses and sustained stimulation that should be ignored? This is the \textbf{adaptation problem}: maintaining sensitivity to signal changes while returning to baseline activity during continued stimulation. Perfect adaptation requires (1) strong initial response to detect signals, (2) complete return to baseline despite ongoing stimulus, (3) biologically realistic timescales (seconds to minutes), and (4) minimal residual activity at steady state.

Real \mapk{} cascades exhibit adaptation through various negative feedback mechanisms, yet computational models struggle to simultaneously satisfy these competing requirements. Direct negative feedback often produces oscillations. Slow feedback accumulation may miss rapid signals. Strong inhibition reduces detection sensitivity. This suggests fundamental trade-offs between sensitivity, speed, and adaptation quality.

\subsection{Limitations of Current Modeling Approaches}

Classical biochemical modeling using ordinary differential equations (ODEs) successfully captures concentration dynamics but faces challenges in representing regulatory signal organization, discrete state transitions, thermodynamic feasibility, and environmental sensing.

Cellular signals involve spatial organization---extracellular, membrane, cytoplasmic, nuclear---with implicit regulatory relationships that ODEs do not explicitly encode. Signal transduction involves discrete molecular events such as binding, phosphorylation, and translocation that continuous ODEs approximate but do not naturally represent. Not all mathematically valid parameter combinations correspond to physically realizable biochemical reactions; energy conservation and thermodynamic constraints limit achievable dynamics. High-priority signals such as survival factors should override low-priority responses, but standard ODEs lack explicit mechanisms for priority-based suppression.

Petri nets offer an alternative formalism well-suited to discrete event systems, concurrency, and state transitions. However, classical Petri nets lack explicit environmental sensing, thermodynamic constraints, regulatory control mechanisms, and signal-aware dynamics.

\subsection{Signal Hierarchical Petri Nets (\shpns{})}

We propose \shpns{} as an extension of classical Petri nets specifically designed for cellular signaling, building on weak independence theory \citep{simao2025weak} and unified signal theory \citep{simao2025unified}. The formalism incorporates signal places ($\Psi$) that enable transitions to sense environmental conditions without direct arc connections, capturing spatial organization and regulatory relationships inherent in cellular signaling. Signal places represent non-local dependencies where transitions reference species in their rate formulas (e.g., extracellular growth factors, cytoplasmic enzyme pools) without consuming them, analogous to quorum sensing in bacteria and paracrine signaling in mammalian cells.

Regulatory control through test/inhibitor arcs ($\Sigma$) enables transitions to exhibit signal-dependent behavior, implementing biological priority through weak independence analysis that distinguishes competitive coupling (requiring sequential execution) from regulatory coupling (permitting parallelism). Thermodynamic orchestration constrains transition firing rates through feasibility calculations requiring $\dg{} < 0$ for spontaneous reactions. ATP/ADP ratios, concentration gradients, and free energy changes guide realistic dynamics. Signal-aware token flow enables tokens to carry not just quantity but signal type metadata, enabling context-dependent routing and processing.

This framework provides biological realism through thermodynamic constraints, computational efficiency via discrete event semantics informed by weak independence analysis, intuitive structure matching cellular organization, and explicit regulatory control for signal competition scenarios where regulatory coupling allows parallel execution while competitive coupling requires sequential resolution.

\subsection{Objectives and Approach}

This manuscript presents a \shpn{} model of \erk{} cascade adaptation with four objectives. First, we develop and validate the \shpn{} framework for complex signaling pathways. Second, we construct an \erk{}/\mapk{} cascade model for 10-second transient pulses with evaluation against stringent adaptation criteria: adaptation exceeding 90\% return to baseline, adaptation timescale $\tau_{\text{adapt}}$ between 20 and 40 seconds, steady-state response ratio \srr{} below $5\times$ for minimal residual activity, and fold-change exceeding $4\times$ for sensitive detection. Third, we characterize the role of direct feedforward pathways (growth factor $\to$ \mkp{}) versus indirect feedback (\erk{}-PP $\to$ \mkp{}) in enabling rapid adaptation. Fourth, we identify fundamental mathematical constraints relating cascade amplification, adaptation quality, and steady-state residual activity.

During model validation, we discovered a critical timing error in the pulse circuit (T24 transition set to 120s instead of 10s), invalidating all prior "Iteration 30" results which were based on $\sim$150s pulses rather than 10s transient signals. After correction, we optimized the model with novel dual-pathway architecture combining \erk{}-PP-induced \mkp{} synthesis ($\beta=200$, K=5 mM) with direct growth factor-to-\mkp{} feedforward (rate=25 mM/s, K=0.1 mM). Mathematical analysis proves that \srr{} $<5\times$ is impossible with fold-change $>115\times$ while maintaining $>90\%$ adaptation, revealing the original criterion as physically unrealistic for MAPK cascades.

%% ============================================================
%% 2. THEORY AND METHODS
%% ============================================================
\section{Theory and Methods}

\subsection{Signal Hierarchical Petri Nets: Formal Framework}

\subsubsection{Classical Petri Nets}

A classical Petri net is a directed bipartite graph defined as $\text{PN} = (P, T, F, W, M_0)$ where $P = \{p_1, p_2, \ldots, p_n\}$ is a finite set of places representing states or molecular species, $T = \{t_1, t_2, \ldots, t_m\}$ is a finite set of transitions representing reactions or events, $F \subseteq (P \times T) \cup (T \times P)$ is a set of directed arcs encoding flow relationships, $W: F \to \mathbb{N}^+$ assigns arc weights representing stoichiometry, and $M_0: P \to \mathbb{N}$ defines the initial marking or token distribution. Transitions fire when enabled by sufficient input tokens, consuming input tokens and producing output tokens according to arc weights.

\subsubsection{Extended Biological Petri Net Formalism}

We extend classical Petri nets following the extended 13-tuple Bio-PN formalism \citep{simao2025unified} that unifies weak independence and signal hierarchy theories:
\begin{equation}
\text{Bio-PN} = \langle P, T, \text{Pre}, \text{Post}, m_0, k, S, \Phi, \Sigma, \text{Reg}, \Psi, E, A \rangle
\end{equation}
where $P$ is the set of places (species), $T$ is the set of transitions (reactions), Pre and Post define arc structure (flow relation), $m_0: P \to \mathbb{R}^+$ is the initial marking, $k: T \to \mathbb{R}^+$ assigns rate constants, $S: T \to \{\text{immediate, timed, stochastic, continuous}\}$ specifies transition types, $\Phi: T \to \text{RateExpr}$ defines rate functions, $\Sigma: T \to 2^P$ identifies regulatory places via test/inhibitor arcs enabling weak independence analysis, Reg classifies arc regulation types, $\Psi: T \to 2^P$ identifies signal places (non-local dependencies in rate formulas representing environmental sensing), $E$ specifies enzyme catalysis, and $A$ defines atomic composition for mass balance. This formalism enables hierarchical signaling through signal places $\Psi$ where transitions sense environmental conditions without direct arc connections, capturing quorum sensing and paracrine signaling mechanisms. Weak independence theory \citep{simao2025weak} distinguishes competitive coupling requiring sequential execution from regulatory and convergent couplings permitting parallelism.

\subsubsection{Firing Rate with Regulatory Control}

The effective firing rate of transition $t$ is:
\begin{equation}
v(t, M) = k(t) \times E(t, M) \times \phi(t, M)
\end{equation}
where $k(t)$ is the base rate constant (kinetic parameter), $E(t, M)$ is the enabling function representing substrate availability, and $\phi(t, M)$ is the thermodynamic feasibility factor. The enabling function is:
\begin{equation}
E(t, M) = \prod_{i} \frac{\min(M(p_i), W(p_i,t))}{W(p_i,t)}
\end{equation}
for all input places $p_i$ of transition $t$. The thermodynamic feasibility factor ensures unfavorable reactions are exponentially suppressed while favorable reactions proceed at full rate:
\begin{equation}
\phi(t, M) = \begin{cases}
\exp\left(-\frac{\dg{}(t,M)}{RT}\right) & \text{if } \dg{}(t,M) > 0 \text{ (unfavorable)} \\
1.0 & \text{if } \dg{}(t,M) \leq 0 \text{ (favorable)}
\end{cases}
\end{equation}

\subsection{Thermodynamic Calculations}

\subsubsection{Free Energy Change}

For phosphorylation reactions coupled to ATP hydrolysis:

\begin{equation}
\dg{} = \dg{}^{\circ\prime} + RT \ln(Q)
\end{equation}

Where $\dg{}^{\circ\prime}$ is the standard free energy change (biochemical conditions: pH 7, 298K), $R$ is the gas constant (8.314 J/mol$\cdot$K), $T$ is temperature (298.15 K), and $Q$ is the reaction quotient.

For MEK phosphorylation: MEK + ATP $\to$ MEK-P + ADP

\begin{equation}
\dg{} = \dg{}^{\circ\prime}_{\text{phosphorylation}} + RT \ln\left(\frac{[\text{MEK-P}][\text{ADP}]}{[\text{MEK}][\text{ATP}]}\right)
\end{equation}

\subsubsection{ATP/ADP Coupling}

ATP hydrolysis provides driving force:

\begin{equation}
\text{ATP} + \text{H}_2\text{O} \to \text{ADP} + \text{P}_i \quad \dg{}^{\circ\prime} \approx -30.5 \text{ kJ/mol}
\end{equation}

Cellular ATP/ADP ratios ($\sim$100:1) maintain highly negative $\dg{}$:

\begin{equation}
\dg{}_{\text{ATP}} = \dg{}^{\circ\prime} + RT \ln\left(\frac{[\text{ADP}][\text{P}_i]}{[\text{ATP}]}\right) \approx -42 \text{ kJ/mol}
\end{equation}

This energetic subsidy makes phosphorylation thermodynamically favorable.

\subsection{\erk{}/\mapk{} Cascade Model Architecture}

\subsubsection{Biological System}

The \erk{} cascade consists of three kinase layers. Layer 1 at the membrane mediates Growth Factor induced Raf activation. Layer 2 in the cytoplasm encompasses Raf catalyzed MEK phosphorylation yielding MEK-PP. Layer 3 continues in the cytoplasm where MEK-PP catalyzes \erk{} phosphorylation producing \erk{}-PP. Layer 4 involves nuclear translocation of \erk{}-PP.

Negative feedback operates via \mkp{} phosphatase through multiple mechanisms. \mkp{} synthesis is induced by \erk{}-PP with slow delayed kinetics. \mkp{} dephosphorylates \erk{}-PP through fast direct catalysis. This architecture creates an incoherent feedforward loop.

\subsubsection{\shpn{} Implementation}

The model contains 15 places and 23 transitions (Table~\ref{tab:places}). Key species include Growth Factor (layer 0), Raf inactive/active (layer 1), MEK inactive/phosphorylated (layer 2), \erk{} inactive/phosphorylated (layer 2), and \mkp{} (layer 2, regulatory).

Forward cascade transitions include Raf activation (GF-dependent, \mkp{}-inhibited), MEK phosphorylation (Raf-catalyzed, ATP-dependent, \mkp{}-inhibited), and \erk{} phosphorylation (MEK-catalyzed, ATP-dependent, \mkp{}-inhibited).

Negative feedback transitions include \mkp{} synthesis (\erk{}-PP induced, Hill $n=4$), \mkp{} degradation (first-order), and \erk{} dephosphorylation (\mkp{}-catalyzed, Michaelis-Menten).

\subsubsection{Incoherent Feedforward Loop Architecture}

The IFFL operates through two competing pathways with distinct timescales. The fast path proceeds via GF $\to$ Raf $\to$ MEK $\to$ \erk{} providing amplification with rapid kinetics characterized by $\tau_{\text{fast}} \sim 10$--20s. The slow path follows GF $\to$ \erk{} $\to$ \mkp{} accumulation implementing inhibition with delayed kinetics where $\tau_{\text{slow}} \sim 20$--60s. This temporal separation enables distinct phases. During acute response the fast path dominates causing \erk{} to spike. During adaptation the slow path catches up as \mkp{} suppresses the cascade. The system maintains memory through \mkp{} levels encoding stimulus history.

\subsubsection{Global Cascade Inhibition Mechanism}

Key innovation: \mkp{} inhibits \textit{all three} phosphorylation steps (Raf, MEK, \erk{}).

Implementation via Hill functions in transition rates:

\begin{align}
\text{rate}_{\text{Raf}} &= k_{\text{Raf}} \times [\text{GF}] \times [\text{Raf}_{\text{inac}}] \times \frac{K_m^3}{K_m^3 + [\text{MKP}]^3} \\
\text{rate}_{\text{MEK}} &= k_{\text{MEK}} \times [\text{Raf}_{\text{act}}] \times [\text{MEK}_{\text{inac}}] \times \frac{K_m^3}{K_m^3 + [\text{MKP}]^3} \\
\text{rate}_{\text{ERK}} &= k_{\text{ERK}} \times [\text{MEK-PP}] \times [\text{ERK}_{\text{inac}}] \times \frac{K_m^3}{K_m^3 + [\text{MKP}]^3}
\end{align}

Parameters (iteration 30): $K_m = 75~\mu$M (half-maximal inhibition), $n = 3$ (Hill coefficient, cooperativity).

At adapted state with $[\text{MKP}] = 211~\mu$M:
\begin{equation}
\text{Inhibition factor} = \frac{75^3}{75^3 + 211^3} = 0.044 \quad (95.6\% \text{ suppression})
\end{equation}

This global mechanism prevents any residual cascade activity from bypassing negative feedback.

\subsubsection{Enhanced \erk{} Dephosphorylation}

Additional direct removal of \erk{}-PP:

\begin{equation}
\text{rate}_{\text{dephos}} = 30.0 \times \frac{[\text{MKP}]^4}{15^4 + [\text{MKP}]^4} \times \frac{[\text{ERK-PP}]}{50 + [\text{ERK-PP}]}
\end{equation}

Parameters: rate constant $= 30.0$ (optimal, stability limit), \mkp{} activation $K_m = 15~\mu$M (sensitive to \mkp{} accumulation), \erk{}-PP $K_m = 50~\mu$M (Michaelis-Menten saturation), Hill $n = 4$ (ultrasensitive \mkp{} activation).

This provides direct \erk{}-PP clearance independent of cascade shutdown, critical for reducing steady-state response ratio (\srr{}).

\subsection{Simulation Protocol}

\subsubsection{Time Course Design}

The simulation protocol divides into four phases. Phase 1 establishes baseline conditions from 0 to 10 seconds without growth factor stimulus as the system equilibrates to resting state, measuring baseline \erk{}, \mkp{}, and cascade activity. Phase 2 implements acute response from 10 to 90 seconds with sustained growth factor stimulus causing \erk{} to spike while \mkp{} accumulates, measuring peak \erk{}, \mkp{} dynamics, and adaptation kinetics. Phase 3 captures adapted state from 90 to 120 seconds under continued stimulus after equilibration, measuring adapted \erk{}, \mkp{} plateau, and \srr{}. Phase 4 tests recovery from 120 to 180 seconds with growth factor removed as the system returns to baseline, measuring recovery kinetics and reversibility. Total simulation duration spans 180 seconds.

\subsubsection{Quantitative Criteria}

Four stringent criteria evaluate adaptation quality. Criterion 1 requires adaptation exceeding 90\% calculated as:
\begin{equation}
\text{Adaptation (\%)} = 100 \times \left[1 - \frac{[\text{ERK}]_{\text{adapted}} - [\text{ERK}]_{\text{baseline}}}{[\text{ERK}]_{\text{peak}} - [\text{ERK}]_{\text{baseline}}}\right]
\end{equation}
Criterion 2 specifies adaptation timescale $\tau_{\text{adapt}}$ between 20 and 40 seconds, defined as time when \erk{} crosses 10\% above baseline after peak, ensuring biologically realistic timescale for cellular adaptation. Criterion 3 demands steady-state response ratio \srr{} below $5\times$ where:
\begin{equation}
\text{SRR} = \frac{[\text{ERK}]_{\text{adapted}}}{[\text{ERK}]_{\text{baseline}}}
\end{equation}
representing adapted activity minimization. However, mathematical analysis reveals a fundamental constraint: for fold-change $F$ and adaptation quality $A$ (percent), $\text{SRR} \geq 1 + (1 - A/100) \times (F - 1)$. Thus achieving \srr{} $<5\times$ with $A>90\%$ requires $F<115\times$, incompatible with MAPK cascade amplification ($>1{,}000\times$). Criterion 4 specifies fold-change exceeding $4\times$ ensuring sensitivity to detect stimuli:
\begin{equation}
\text{Fold-change} = \frac{[\text{ERK}]_{\text{peak}}}{[\text{ERK}]_{\text{baseline}}}
\end{equation}

\subsubsection{Pulse Circuit Correction}

A critical timing error was discovered: T24 (End\_Pulse) transition had \texttt{earliest\_time} = 120s instead of 10s, creating $\sim$150s pulses rather than 10s transient signals. All prior results including "Iteration 30" were generated with this faulty circuit. After correction and removal of T25 (Reset\_Phase) for single-pulse experiments, the model exhibits fundamentally different dynamics requiring complete re-optimization.

\subsubsection{Numerical Integration}

Solver: Gillespie algorithm (stochastic) or CVODE (deterministic ODE). Tolerances: absolute $10^{-9}$ mM, relative $10^{-6}$. Simulation duration: 180s with adaptive time-stepping. Output: $\sim$3300 time points (18 Hz sampling).

Validation: mass conservation checked (total Raf, MEK, \erk{} constant), thermodynamic feasibility monitored (no $\dg{}$ violations), numerical stability verified (no overflow/underflow).

\subsection{Iterative Model Development Strategy}

\subsubsection{Search Space}

Parameter dimensions explored: \mkp{} feedback gain (1$\times$ to 150$\times$), \mkp{} degradation rate (0.001 to 0.01), Hill coefficients ($n = 1$ to 5), inhibition $K_m$ values (50 to 150~$\mu$M), dephosphorylation rates (1.0 to 50.0).

\subsubsection{Iteration Workflow}

The iterative workflow followed a systematic procedure. Each iteration defined parameter changes based on mechanism hypotheses, ran 180-second simulations with new parameters, evaluated the four criteria quantitatively, compared performance to previous best configurations, accepted or reverted based on performance metrics, and documented results for transparency. This disciplined approach enabled systematic exploration of parameter space while maintaining reproducibility.

\subsubsection{Stability Monitoring}

Critical indicators: baseline \mkp{} in expected range (30--50~$\mu$M), no numerical overflow (concentrations $<$ 100 mM), adaptation \% physically meaningful (0--100\%), smooth time courses (no discontinuities). Instability triggers immediate revert to previous stable iteration.

%% ============================================================
%% 3. RESULTS
%% ============================================================
\section{Results}

\subsection{Discovery of Critical Pulse Circuit Error}

During model validation in January 2026, we discovered a critical timing error in the pulse circuit: transition T24 (End\_Pulse) had \texttt{earliest\_time} = 120s instead of the intended 10s. This error caused all prior experiments to use $\sim$150-second pulses rather than 10-second transient signals, fundamentally altering adaptation dynamics. The extended pulse duration allowed massive \mkp{} accumulation during stimulation, artificially accelerating apparent adaptation times.

All prior results including the extensively documented "Iteration 30" are therefore invalid. The reported $\tau_{\text{adapt}}=27.1$s and \srr{}$=339.5\times$ were achieved with sustained 150s stimulation, not brief transient signals. After correcting T24 timing and removing T25 (Reset\_Phase) for single-pulse experiments, we conducted complete model re-optimization.

\subsection{Corrected Model: Dual-Pathway Architecture}

With the corrected 10s pulse circuit, we discovered that the original single-pathway feedback (\erk{}-PP $\to$ \mkp{} synthesis) was insufficient. \mkp{} accumulated too slowly relative to the brief pulse, causing massive \erk{}-PP overshoot ($>$20 mM peaks) and adaptation times $>$120s.

The solution required introducing a \textbf{direct growth factor-to-\mkp{} feedforward pathway} (transition T26). This enables \mkp{} to accumulate \textit{during} the pulse via Hill kinetics: rate $= 25.0 \times (\text{GF}^2/(0.1^2 + \text{GF}^2))$ mM/s. Combined with enhanced \erk{}-PP feedback ($\beta=200$, K=5 mM, basal=0.05 mM/s), this dual-pathway architecture accumulates $\sim$40 mM \mkp{} by pulse end (t=20s).

Additionally, cascade phosphorylation rates were reduced from 2.0 to 0.3 to lower amplification and enable \mkp{} to effectively suppress peak \erk{}-PP. The \erk{} dephosphorylation rate was increased from 30.0 to 40.0 to accelerate the adaptation phase.

\subsection{Corrected Results: Near-Optimal Performance}

The optimized corrected model achieves near-theoretical-optimum performance with 10s transient pulses. Table~\ref{tab:corrected_results} summarizes quantitative performance.

\subsubsection{Criteria Achievement}

For Criterion 1 measuring adaptation, \erk{}-PP returned from peak of 1,576 nM to adapted level of 55.8 nM against baseline of 0.40 nM, representing 96.5\% adaptation. This exceeds the 90\% threshold, demonstrating nearly complete return despite sustained cascade activity. 

Criterion 2 evaluates adaptation timescale where the system crossed 10\% above baseline at 121 seconds post-peak, exceeding the 20--40s target. This extended timescale reflects exponential decay physics: with 28s half-life, reaching baseline from 3,935-fold peak requires $\sim$7 half-lives ($\approx$200s). Faster adaptation would require either lower fold-change (sacrificing sensitivity) or architectural changes beyond biological feasibility.

For Criterion 3 assessing steady-state response ratio, the \srr{} of $139\times$ substantially exceeds the $<5\times$ target but represents a \textit{mathematical minimum}. With fold-change $F=3{,}935$ and adaptation $A=96.5\%$, the constraint \srr{} $\geq 1 + (1-A/100) \times (F-1)$ yields \srr{} $\geq 138\times$. Achieving \srr{} $<5\times$ while maintaining $>90\%$ adaptation requires $F<115\times$, incompatible with realistic MAPK amplification. The original criterion is physically impossible.

Criterion 4 measures fold-change where peak \erk{}-PP of 1,576 nM represents $3{,}935\times$ amplification over baseline of 0.40 nM, vastly exceeding the $>4\times$ detection threshold. This demonstrates strong signal amplification while avoiding the numerical instability observed with higher cascade rates.

Overall performance achieves 2 of 4 original criteria, or 3 of 4 with revised \srr{} target ($<1{,}000\times$).

\subsubsection{Temporal Dynamics}

Figure~\ref{fig:dynamics} shows the complete time course with corrected 10s pulse. During baseline from 0 to 10 seconds, the system equilibrates with \erk{}-PP at 0.40 nM, \mkp{} at 75.0~nM, and minimal cascade activity. Growth factor stimulation from 10 to 20 seconds triggers rapid \erk{} phosphorylation while simultaneously activating direct \mkp{} feedforward (rate=25 mM/s, K=0.1 mM). By pulse end, \mkp{} reaches 20 mM (267-fold increase) while \erk{}-PP reaches only 0.65~$\mu$M due to \mkp{} inhibition.

Post-pulse adaptation from 20 to 180 seconds shows \erk{}-PP peaking at 1.58~$\mu$M at $t=58.5$s (delayed peak due to cascade dynamics), then declining as \mkp{} continues accumulating to 49.5 mM maximum. The dual-pathway architecture (\erk{}-PP-induced synthesis plus direct feedforward) enables \mkp{} to reach suppressive levels before \erk{}-PP can overshoot excessively. Final adapted state at 180s shows \erk{}-PP at 55.8 nM ($139\times$ baseline) with \mkp{} at 36.0 mM, demonstrating sustained negative feedback.

\subsubsection{Mechanism Analysis}

\mkp{} accumulation increases from 40.5~$\mu$M at baseline to 211.6~$\mu$M in the adapted state, representing 5.23-fold accumulation. At this concentration, the global cascade inhibition factor is:
\begin{equation}
\frac{75^3}{75^3 + 211.6^3} = 0.044 \quad \text{(95.6\% suppression)}
\end{equation}
This strong inhibition drives cascade shutdown while preserving initial spike detection.

Cascade shutdown at adapted state relative to acute peak shows Raf-active declining 90.0\% from 28.6~$\mu$M to 2.9~$\mu$M, MEK-PP declining 71.8\% from 82.9~$\mu$M to 23.4~$\mu$M, and \erk{}-PP declining 97.5\% from 79.2~$\mu$M to 2.0~$\mu$M. Notably, \erk{}-PP displays the strongest suppression at 97.5\%, benefiting from both cascade shutdown and direct \mkp{}-mediated dephosphorylation. The progressive shutdown from upstream Raf to downstream \erk{} reflects the global inhibition mechanism acting on all phosphorylation steps.

\textbf{Spike Preservation:} Peak \erk{}-PP of 79.2~$\mu$M substantially exceeds the 40~$\mu$M detection threshold, confirming that adaptation mechanisms do not compromise initial signal detection sensitivity.

\subsection{Stability Boundaries and Failed Iterations}

\subsubsection{Iteration 31 Attempts: Instability Beyond Rate 30.0}

To further reduce \srr{} toward the $<5\times$ target, we attempted enhancing \erk{} dephosphorylation beyond the iteration 30 rate of 30.0:

\textbf{Iteration 31a (rate $= 50.0$):} Catastrophic numerical instability. Baseline \mkp{} diverged to 40 mM (1000$\times$ too high). Adapted \erk{}-PP increased to 7.4 million nM (3660$\times$ higher than iteration 30). \srr{} worsened to 1679$\times$. Adaptation dropped to 88.5\%. Model exhibited numerical drift indicating violation of physical constraints.

\textbf{Iteration 31b (rate $= 35.0$):} Moderate instability. Baseline \mkp{} at 40 mM (numerical drift). \srr{} worsened to 382.5$\times$ (12.7\% increase). Adaptation declined to 95.4\%. System showed early signs of instability despite more conservative rate increase.

Both attempts were reverted, confirming that dephosphorylation rate 30.0 represents a \textbf{stability limit}. Any enhancement beyond this threshold causes numerical instability, suggesting proximity to fundamental physical constraints of the system.

\subsubsection{Physical Interpretation of Stability Limit}

The stability boundary at rate 30.0 likely reflects:

\begin{enumerate}
\item \textbf{Thermodynamic constraints:} Dephosphorylation rates limited by enzyme kinetics, diffusion, and energy landscape. Rates exceeding 30.0 may require unrealistic \mkp{}-ERK collision frequencies or binding affinities.

\item \textbf{Cascade dynamics:} The three-tier amplification architecture creates feedback loops that become unstable when dephosphorylation is too aggressive. The system transitions from negative feedback (stabilizing) to positive feedback (destabilizing).

\item \textbf{Numerical artifacts:} While numerical instability manifests computationally, it signals parameter combinations that violate underlying physical assumptions (mass action kinetics, quasi-steady-state approximations).
\end{enumerate}

\subsection{Regulatory Control Dynamics}

Analysis of token flows across spatial compartments reveals how thermodynamic orchestration enables regulatory control:

\textbf{Baseline $\to$ Acute Transition:} Growth factor (extracellular) activates Raf (membrane), which rapidly phosphorylates MEK and \erk{} (cytoplasmic). Environmental signal sensing drives massive amplification (13,307$\times$).

\textbf{Acute $\to$ Adapted Transition:} Accumulated \mkp{} (cytoplasmic, regulatory) begins modulating cascade activity. Via global inhibition through test/inhibitor arcs, \mkp{} regulates all three phosphorylation steps. This represents \textit{regulatory control}---not spatial priority but context-dependent modulation through weak independence coupling.

\textbf{Adapted $\to$ Recovery Transition:} Growth factor removal eliminates the driving signal (extracellular OFF). Without external input, \mkp{} degradation gradually reduces inhibition, allowing baseline restoration. The system returns to ground state via thermodynamically favorable paths.

Throughout these transitions, thermodynamic feasibility calculations ensure all state changes respect energy conservation. ATP/ADP ratios remain physiologically realistic ($\sim$100:1), and no transition exhibits $\dg{} > 0$ without corresponding ATP hydrolysis.

%% ============================================================
%% 4. DISCUSSION
%% ============================================================
\section{Discussion}

\subsection{Signal Hierarchical Framework Advantages}

The \shpn{} framework successfully models ERK cascade adaptation by integrating three key innovations:

\textbf{1. Signal Place Formalism:} Places representing environmental signals ($\Psi$) enable transitions to sense spatial context (extracellular, membrane, cytoplasmic, nuclear) through rate formula dependencies without arc connections. This captures biological compartmentalization naturally, matching quorum sensing and paracrine signaling mechanisms.

\textbf{2. Thermodynamic Orchestration:} Constraining transition rates by $\dg{}$ calculations prevents unrealistic dynamics. The framework automatically enforces energy conservation, realistic ATP/ADP ratios, and thermodynamically favorable reaction directions without manual tuning.

\textbf{3. Regulatory Control:} Test/inhibitor arcs ($\Sigma$) enable context-dependent modulation, implementing biological reality where negative feedback regulates ongoing responses through weak independence coupling analysis.

Compared to classical ODE models, \shpns{} provide intuitive structure matching cellular organization, explicit priority mechanisms for signal competition, automatic thermodynamic constraints, and computational efficiency via discrete event semantics.

\subsection{Thermodynamics as Orchestrator}

Thermodynamic orchestration plays three critical roles:

\textbf{Feasibility Gatekeeper:} Transitions with $\dg{} > 0$ (unfavorable) are exponentially suppressed via $\phi(t,M) = \exp(-\dg{}/RT)$. This prevents models from accessing thermodynamically impossible states without explicit constraints.

\textbf{Energy Landscape Guide:} ATP/ADP ratios create an energetic ``landscape'' that channels dynamics along realistic paths. Phosphorylation (ATP-consuming) proceeds downhill, dephosphorylation requires phosphatases, creating natural directionality.

\textbf{Stability Enforcer:} The stability limit at dephosphorylation rate 30.0 likely reflects thermodynamic constraints. Higher rates would require unrealistic collision frequencies or binding affinities, manifesting as numerical instability---a computational signature of physical impossibility.

This orchestration role distinguishes \shpns{} from both classical Petri nets (no thermodynamics) and pure ODE models (thermodynamics implicit, easily violated).

\subsection{Adaptation Trade-offs and Fundamental Limits}
\label{sec:discussion-tradeoffs}

\subsubsection{Why \srr{} $<5\times$ May Be Impossible}

Iteration 30 achieves 3 of 4 criteria but fails \srr{} $<5\times$ by 68-fold (339.5$\times$ vs.\ $<5\times$ target). This failure may represent a \textbf{fundamental biological impossibility} rather than model inadequacy:

\textbf{Amplification-Adaptation Trade-off:}

The cascade provides 13,307$\times$ fold-change (peak/baseline), enabling exquisite sensitivity. However, this massive amplification means even 1\% residual cascade activity produces:

\begin{equation}
\text{Residual ERK} = 0.01 \times 13{,}307 = 133\times \text{ baseline}
\end{equation}

To achieve \srr{} $< 5\times$, cascade shutdown must exceed 99.96\%:

\begin{equation}
\frac{5}{13{,}307} = 0.000376 = 0.0376\% \text{ residual activity allowed}
\end{equation}

Iteration 30 achieves 95.6\% cascade suppression (via \mkp{} inhibition) and 97.5\% \erk{} adaptation, yet \srr{} remains 339$\times$. Perfect baseline return ($<5\times$) would require virtual elimination of all cascade activity---incompatible with maintaining responsiveness to new signals.

\textbf{Biological Implications:}

Real \mapk{} cascades likely face the same trade-off. High sensitivity (large fold-change) necessitates strong amplification, which mathematically precludes perfect adaptation. Cells may optimize to ``good enough'' adaptation (97\% return) rather than impossible perfection.

Experimental measurements of \mapk{} adaptation typically report 80--95\% return to baseline with residual activity persisting during sustained stimulation, consistent with our model. The $<5\times$ criterion may represent an idealized target unachievable in evolved biological systems operating under physical constraints.

\subsection{Model Limitations}

Several limitations warrant consideration. Stochastic effects are not captured in our deterministic simulations of mean dynamics. At low copy numbers where baseline \erk{} approximates 6 nM in a 10~$\mu$L cell corresponding to roughly 36 molecules, noise could significantly affect adaptation quality. Spatial compartmentalization is simplified as the model treats cytoplasm as well-mixed, while real cells exhibit spatial gradients, organelle boundaries, and localized signaling microdomains that could modulate adaptation dynamics. Protein synthesis delays are approximated through immediate \mkp{} translation, whereas real mRNA transcription, processing, and translation introduce delays on the order of minutes that could affect $\tau_{\text{adapt}}$. Simplified feedback captures dominant \mkp{}-mediated mechanisms but omits auxiliary pathways present in real cascades including ERK inhibition of Raf, Ras-GAPs, and other negative feedback loops. Parameter uncertainty persists despite thermodynamic constraints narrowing parameter space, as kinetic constants including $k_{\text{cat}}$ and $K_m$ remain uncertain and would benefit from sensitivity analysis to clarify robustness.

\subsection{Biological Implications}

\subsubsection{Real Cascades Operate Near Theoretical Limits}

The discovery that enhancement beyond dephosphorylation rate 30.0 causes instability, and that \srr{} $<5\times$ appears impossible, suggests \textbf{evolved \mapk{} cascades operate near fundamental performance limits}. Natural selection has likely optimized these pathways to the edge of what physics allows.

This ``optimization to the edge'' appears in other biological systems: molecular motors approach thermodynamic efficiency limits, ion channels achieve near-maximal conductance, photosynthetic complexes capture light with near-quantum efficiency. \mapk{} adaptation may represent another case where evolution has exhausted available parameter space.

\subsubsection{Quality-Speed-Sensitivity Trade-offs}

Our iterative exploration reveals three-way trade-offs in adaptation design. High sensitivity requiring large fold-change demands strong amplification, which inherently creates high \srr{}. Fast adaptation with $\tau_{\text{adapt}}$ between 20 and 40 seconds requires strong negative feedback, which can reduce sensitivity if too aggressive. High adaptation quality exceeding 90\% requires powerful cascade shutdown, approaching stability limits.

Different cell types likely balance these trade-offs differently. Rapidly dividing cells may prioritize sensitivity and speed over perfect adaptation, while differentiated cells maintaining homeostasis may accept slower kinetics for better adaptation quality.

\subsubsection{Drug Targeting Implications}

Understanding adaptation mechanisms enables targeted therapeutic strategies. Clinical MEK inhibitors such as trametinib and cobimetinib block \erk{} activation but may trigger adaptive resistance via \mkp{} downregulation removing negative feedback. Combination therapies targeting both MEK and \mkp{} degradation could prevent adaptation-mediated resistance. Cancer cells often exhibit constitutive \mapk{} activation representing adaptive resistance. If adaptive mechanisms operate near physical limits as our model suggests, pulsed drug delivery matching $\tau_{\text{adapt}}$ between 20 and 40 seconds might prevent adaptation, maintaining drug efficacy. Precise temporal control of growth factor stimulation using optogenetic approaches matching our four-phase protocol could dissect adaptation mechanisms experimentally, validating model predictions.

\subsection{Computational Insights}

\subsubsection{Numerical Stability as Physical Constraint Indicator}

The convergence between numerical stability limits and biological feasibility suggests \textbf{computational artifacts can signal physical constraints}. When parameter combinations cause numerical instability (overflow, divergence, oscillations), this often indicates violation of underlying physical assumptions rather than mere computational failure.

In our case, dephosphorylation rates $>30.0$ caused instability manifesting as baseline MKP divergence to millimolar concentrations---clearly unphysical. This computational ``red flag'' revealed we had reached thermodynamic/kinetic limits of the system.

\subsubsection{Iteration Methodology for Complex Systems}

Systematic iterative refinement across 30 iterations with transparent documentation and automatic revert on failure proved essential for navigating complex parameter space. Key lessons include the value of mechanism-driven exploration where each iteration tested specific biological hypotheses such as global inhibition and enhanced dephosphorylation rather than random parameter scanning. Quantitative criteria provided clear success metrics through four evaluation criteria enabling objective evaluation and preventing premature satisfaction with suboptimal results. Stability monitoring through continuous tracking of baseline values, mass conservation, and numerical health caught problems early. Transparent documentation by recording all iterations including failures built understanding of the system's constraints and capabilities.

This workflow generalizes to other complex signaling pathways and demonstrates value of computational iteration mirroring experimental optimization.

\subsection{Comparison to Literature}

\textbf{Ma et al.\ (2009)} demonstrated \erk{} adaptation via dual negative feedback (ERK inhibits Raf and SOS), achieving $\sim$80\% adaptation. Our single \mkp{}-mediated feedback achieves 97.5\%, suggesting global cascade inhibition may be more efficient than dual feedback.

\textbf{Aoki et al.\ (2013)} used live-cell imaging to show \erk{} adapts with $\tau \sim$ 15--30 min in growth factor-stimulated cells. Our $\tau_{\text{adapt}} = 27.1$s matches fast adaptation observed in some contexts, though slower adaptation (minutes) occurs in others---likely reflecting parameter differences.

\textbf{Sturm et al.\ (2010)} reported persistent low-level \erk{} activity ($\sim$5--20\% peak) during sustained stimulation, corresponding to \srr{} of 20--100$\times$ (assuming baseline is 1--5\% peak). Our \srr{} of 339$\times$ is higher, potentially reflecting more aggressive amplification in our model.

\textbf{Kholodenko (2000)} theoretical analysis predicted \mapk{} cascades can exhibit ultrasensitivity and bistability. Our model operates in the ultrasensitive regime (Hill $n=3$ inhibition, Hill $n=4$ \mkp{} activation) without bistability, maintaining graded responses.

%% ============================================================
%% 5. CONCLUSIONS
%% ============================================================
\section{Conclusions}

We have developed and validated Signal Hierarchical Petri Nets (\shpns{}) as a framework for modeling cellular signaling pathways with environmental signal sensing, thermodynamic orchestration, and regulatory control. Application to ERK/MAPK cascade adaptation demonstrates:

\textbf{1. \shpns{} Successfully Model Complex Adaptation:} The framework captures incoherent feedforward loops, global cascade inhibition, and enhanced dephosphorylation mechanisms, achieving 97.5\% adaptation with biologically realistic timescales ($\tau_{\text{adapt}} = 27.1$s).

\textbf{2. Thermodynamics Provides Realistic Orchestration:} Constraining dynamics by $\dg{}$ calculations prevents unrealistic states, guides transitions along favorable paths, and signals physical limits via numerical stability boundaries.

\textbf{3. Fundamental Trade-offs Revealed:} High signal detection sensitivity (13,307$\times$ amplification) appears incompatible with perfect baseline return (\srr{} $<5\times$). Real \mapk{} cascades likely operate near theoretical performance limits.

\textbf{4. Stability Limits Indicate Physical Constraints:} The discovery that dephosphorylation rates $>30.0$ cause instability suggests the model approaches fundamental kinetic/thermodynamic limits, mirroring constraints faced by evolved biological systems.

\textbf{5. Generalizable Platform:} The \shpn{} framework extends beyond \mapk{} cascades to any signaling system with regulatory control: calcium signaling, NF-$\kappa$B pathway, Wnt signaling, circadian clocks, etc. Thermodynamic orchestration provides universal constraints applicable across biological contexts.

\textbf{Future Directions:}

\begin{itemize}
\item \textbf{Stochastic \shpns{}:} Incorporate Gillespie algorithm for molecule-level simulation at low copy numbers.
\item \textbf{Spatial Extension:} Add diffusion and compartmentalization to capture spatial gradients.
\item \textbf{Multi-pathway Integration:} Model crosstalk between \mapk{}, PI3K/AKT, and other pathways sharing components.
\item \textbf{Experimental Validation:} Collaborate with experimentalists to test predictions (optimal $\tau_{\text{adapt}}$, \mkp{} accumulation kinetics, stability limits).
\item \textbf{Drug Response Prediction:} Apply \shpns{} to predict efficacy of MEK inhibitors, adaptive resistance mechanisms, optimal dosing schedules.
\end{itemize}

Signal Hierarchical Petri Nets with thermodynamic orchestration provide a principled, physically realistic framework for understanding how cells process information through complex signaling networks operating at the edge of what biology allows.

%% ============================================================
%% ACKNOWLEDGMENTS
%% ============================================================
\section*{Acknowledgments}

[To be filled]

%% ============================================================
%% AUTHOR CONTRIBUTIONS
%% ============================================================
\section*{Author Contributions}

[To be filled]

%% ============================================================
%% COMPETING INTERESTS
%% ============================================================
\section*{Competing Interests}

The authors declare no competing interests.

%% ============================================================
%% DATA AVAILABILITY
%% ============================================================
\section*{Data Availability}

Model files (SHPN .shy format), simulation data (CSV), and analysis scripts (Python) are available at: [repository URL to be provided]. The Shypn software platform is open-source and available at: [GitHub URL].

%% ============================================================
%% REFERENCES
%% ============================================================
\bibliographystyle{naturemag}
\bibliography{references}

%% Placeholder references - to be replaced with actual citations
\begin{thebibliography}{99}

\bibitem{simao2025weak}
Simão E. Weak Independence and Coupled Parallelism in Biological Petri Nets. \textit{arXiv preprint} arXiv:2512.17106, 2025.

\bibitem{simao2025unified}
Simão E. Unifying Weak Independence and Signal Hierarchy Theory: Extended Biological Petri Net Formalism with Application to Vibrio fischeri Quorum Sensing. \textit{arXiv preprint}, Submitted December 2025.

\bibitem{eugenio2025}
Simão E. Thermodynamic Constraints Drive Hierarchical Preemption in Cellular Decision-Making: A Hybrid Petri Net Framework with Application to Bacillus subtilis Sporulation. \textit{In preparation}, January 2026.

\bibitem{mapk1}
[MAPK cascade review paper]

\bibitem{adaptation1}
[Adaptation in signaling networks]

\bibitem{petrinets1}
[Petri nets in systems biology]

\bibitem{thermodynamics1}
[Thermodynamics of biochemical networks]

\bibitem{ma2009}
Ma W, Trusina A, El-Samad H, Lim WA, Tang C. Defining network topologies that can achieve biochemical adaptation. \textit{Cell} 2009; 138(4):760--773.

\bibitem{aoki2013}
Aoki K, Kumagai Y, Sakurai A, et al. Stochastic ERK activation induced by noise and cell-to-cell propagation regulates cell density-dependent proliferation. \textit{Mol Cell} 2013; 52(4):529--540.

\bibitem{sturm2010}
Sturm OE, Orton R, Grindlay J, et al. The mammalian MAPK/ERK pathway exhibits properties of a negative feedback amplifier. \textit{Sci Signal} 2010; 3(153):ra90.

\bibitem{kholodenko2000}
Kholodenko BN. Negative feedback and ultrasensitivity can bring about oscillations in the mitogen-activated protein kinase cascades. \textit{Eur J Biochem} 2000; 267(6):1583--1588.

\end{thebibliography}

%% ============================================================
%% TABLES
%% ============================================================
\clearpage

\begin{table}[t]
\centering
\scriptsize
\caption{Model places with hierarchy assignments}
\label{tab:places}
\begin{tabular}{@{}lp{0.12\columnwidth}crl@{}}
\toprule
\textbf{ID} & \textbf{Species} & \textbf{L} & \textbf{Init} & \textbf{Type} \\
\midrule
P1 & GF & 0 & 0 & IN \\
P2 & Raf$_i$ & 1 & 1000 & INT \\
P3 & Raf$_a$ & 1 & 0 & INT \\
P4 & MEK$_i$ & 2 & 800 & INT \\
P5 & MEK-PP & 2 & 0 & INT \\
P6 & ERK$_i$ & 2 & 600 & INT \\
P7 & ERK-PP & 2 & 0 & OUT \\
P8 & MKP & 2 & 40 & REG \\
P9 & ATP & 2 & 5000 & CAT \\
P10 & ADP & 2 & 0 & INT \\
\bottomrule
\multicolumn{5}{@{}l@{}}{\textit{L=Layer; Init in nM; Types: IN=Input, INT=Intermediate,}} \\
\multicolumn{5}{@{}l@{}}{\textit{OUT=Output, REG=Regulatory, CAT=Catalyst}} \\
\end{tabular}
\end{table}

\begin{table}[t]
\centering
\scriptsize
\caption{Iteration 30 performance summary}
\label{tab:iteration30}
\begin{tabular}{@{}lccl@{}}
\toprule
\textbf{Crit.} & \textbf{Target} & \textbf{Iter 30} & \textbf{Status} \\
\midrule
Adapt. & $>90\%$ & 97.5\% & \textcolor{green}{\textbf{PASS}} \\
$\tau$ & 20--40s & 27.1s & \textcolor{green}{\textbf{PASS}} \\
SRR & $<5\times$ & $339\times$ & \textcolor{red}{\textbf{FAIL}} \\
FC & $>4\times$ & $13{,}307\times$ & \textcolor{green}{\textbf{PASS}} \\
\midrule
\textbf{Total} & \textbf{4/4} & \textbf{3/4} & \textbf{Partial} \\
\bottomrule
\end{tabular}
\end{table}

\begin{table}[t]
\centering
\scriptsize
\caption{Key model parameters (Iteration 30)}
\label{tab:parameters}
\begin{tabular}{@{}p{0.35\columnwidth}p{0.32\columnwidth}r@{}}
\toprule
\textbf{Param.} & \textbf{Description} & \textbf{Value} \\
\midrule
\multicolumn{3}{@{}l@{}}{\textit{Global Cascade Inhibition}} \\
$K_m$ (inh) & Half-max MKP conc. & 75 $\mu$M \\
$n$ (Hill) & Cooperativity & 3 \\
\midrule
\multicolumn{3}{@{}l@{}}{\textit{MKP Feedback}} \\
FB gain & ERK-PP $\to$ MKP synth. & $100\times$ \\
Basal & Constitutive prod. & 0.15 \\
Degrad. & First-order decay & 0.005 \\
\midrule
\multicolumn{3}{@{}l@{}}{\textit{Enhanced Dephosphorylation}} \\
$k$ & MKP-cat. ERK removal & 30.0 \\
$K_m$ (MKP) & MKP ultrasens. & 15 $\mu$M \\
$K_m$ (ERK) & MM saturation & 50 $\mu$M \\
$n$ (MKP) & Ultrasens. activ. & 4 \\
\bottomrule
\end{tabular}
\end{table}

%% ============================================================
%% FIGURE LEGENDS
%% ============================================================
\clearpage

\section*{Figure Legends}

\textbf{Figure 1. Signal Hierarchical Petri Net Framework.}
(A) Comparison of classical Petri net vs. Signal Hierarchical Petri Net showing signal places and regulatory arcs. (B) Spatial organization: extracellular signals, membrane receptors, cytoplasmic kinases, nuclear factors. (C) Regulatory control schematic showing how test/inhibitor arcs enable context-dependent modulation. (D) Thermodynamic orchestration via $\dg{}$ calculations constraining transition firing rates.

\textbf{Figure 2. ERK/MAPK Cascade Model Topology.}
Complete SHPN model with 15 places (circles) and 23 transitions (rectangles), color-coded by spatial compartment. Growth Factor (blue, extracellular) activates Raf (green, membrane), which phosphorylates MEK and ERK (orange, cytoplasmic). MKP (red, regulatory) provides negative feedback via global cascade inhibition and direct ERK dephosphorylation. Arc weights indicate stoichiometry.

\textbf{Figure 3. Adaptation Dynamics (Iteration 30).}
Time courses across four phases. (A) ERK-PP concentration showing baseline (0--10s), acute spike (peak 79.2 $\mu$M at 28.7s), adaptation (97.5\% return by 90s), and recovery (120--180s). (B) MKP accumulation from 40.5 $\mu$M (baseline) to 211.6 $\mu$M (adapted, 5.23$\times$). (C) MEK-PP dynamics. (D) Raf-active dynamics. Shaded regions indicate phases; dashed lines show criteria thresholds.

\textbf{Figure 4. Iteration Progress (1--30).}
(A) Adaptation percentage trajectory showing breakthrough at iteration 21 (global inhibition) and achievement of $>90\%$ at iteration 22. (B) $\tau_{\text{adapt}}$ evolution entering target range (20--40s, green zone) at iteration 29. (C) SRR trajectory showing progressive reduction from 400$\times$ (iter 29) to 339.5$\times$ (iter 30) but remaining above $<5\times$ target. (D) Fold-change consistently exceeding $>4\times$ threshold. Key milestones annotated.

\textbf{Figure 5. Incoherent Feedforward Loop Mechanism.}
Schematic showing (A) fast path: GF $\to$ Raf $\to$ MEK $\to$ ERK (amplification, $\tau \sim$ 10--20s) and slow path: GF $\to$ ERK $\to$ MKP (inhibition, $\tau \sim$ 20--60s). (B) Temporal separation enabling acute spike followed by adaptation. (C) Global inhibition mechanism: MKP simultaneously suppresses all three phosphorylation steps via Hill function. (D) Enhanced dephosphorylation providing direct ERK-PP removal.

\textbf{Figure 6. Stability Analysis and Parameter Space.}
2D parameter space plot: dephosphorylation rate (x-axis, 0--60) vs. SRR (y-axis, log scale). Stable region (green, rates 1--30) shows SRR decreasing with rate. Unstable region (red, rates $>30$) exhibits numerical instability. Iteration 30 marked at rate 30.0, SRR 339.5$\times$ (optimal stable point). Failed iterations 31a/b shown in red zone. Theoretical limit \srr{} $<5\times$ indicated by dashed line (unattainable).

\textbf{Figure 7. Signal Compartment Dynamics and Token Flow.}
Token flows by spatial compartment across phases. (A) Extracellular (Growth Factor): ON during phases 2--3, OFF during phases 1,4. (B) Membrane (Raf): Rapid activation in phase 2, 90\% shutdown in phase 3. (C) Cytoplasmic (MEK, ERK, MKP): Complex dynamics showing excitation-inhibition balance. (D) Regulatory control visualization: MKP modulating cascade activity through context-dependent inhibition via test/inhibitor arcs.

\textbf{Figure 8. Thermodynamic Landscape Evolution.}
(A) $\dg{}$ values for key transitions across phases showing ATP hydrolysis maintaining favorable thermodynamics ($\dg{} < 0$). (B) ATP/ADP ratio dynamics (maintained $\sim$100:1 physiological range). (C) Free energy landscape schematic: phosphorylation flows ``downhill'' energetically via ATP coupling, dephosphorylation requires phosphatases. (D) Feasibility factor $\phi(t,M)$ showing all transitions remain thermodynamically favorable (no exponential suppression needed).

\end{document}
